\section{The Third Ring of Saturn}

George Santayana is the third of the Saturnine philosophers, following Spinoza and Schopenhauer. Although all three have a similar understanding of the Manifest world, the differ on how to express the Unseen. By dint of human effort, they rise to the top of the Visible world and point their mental telescopes at the sky. There they find the secret sun, which shines so bright, that its light hides all the lesser lights of the Visible World.

Spinoza calls it God. This God has attributes including Thought and Extension, although He cannot be found there. These attributes are his versions of the Names, or even Energies, of God. The multitudinous events we observe in the Manifest world are merely modifications of God, the one and only Substance. Our fulfillment, even salvation, is to learn to love that God.

Schopenhauer maintains that the Manifest world of phenomena is a representation in the human mind, although he doesn't explain why the human state is so privileged to give birth to that world. Ultimate reality cannot be found among those phenomena, but he notices that the will is outside phenomena, that is, it is noumenal. Hence, the Ultimate is pure Will, with no intelligence or purpose. Hence, it cannot be understood. Salvation can only mean to stop producing pain in sentient beings, but, best of all, would be to end existence itself.

\paragraph{The Reluctant Catholic}


Although George Santayana was Spanish --- and retained his citizenship throughout his life --- he was raised in Boston where he got a classical education at the Boston Latin School. He was never comfortable there, so, after receiving an inheritance, he ended up living in Rome during the Fascist era. I know the feeling. After being raised in a section of Somerville that was 75\% Catholic, there was a culture shock when my family relocated to a suburb that was dominated by Protestant descendants of the Puritans.

Santayana was culturally and aesthetically Catholic although he could not bring himself to faith. Thus he saw, in Spinoza, the intellectual alternative. Santayana follows Spinoza in many ways, except he concludes that the Attribute of Extension, i.e., ``matter'', is the ultimate reality. Therefore, Thought must be derivative of matter. However, matter has no I, it cannot be a god.

We see a decline in the three rings of Saturn:

\begin{enumerate}


\item \textbf{Spinoza}: God is the ultimate reality, revealed through Thought and Extension

\item \textbf{Schopenhauer}: Arbitrary Will is the ultimate reality, revealed through our noumenal grasp of the human will

\item \textbf{Santayana}: Matter is the ultimate reality, known intuitively by an ``animal faith''.
\end{enumerate}


\paragraph{The Life of Reason}


\begin{quotex}
Those who cannot remember the past are condemned to repeat it.
\flright{\textit{Reason in Common Sense}}
\end{quotex}


Santayana's first masterpiece was his \emph{Life of Reason}, which addressed the role of reason in Common Sense, Society, Religion, Art, and Science. It is impossible here to summarize such a wide-ranging work. It suffices to point out how expansive the Life of Reason is. Unlike some narrow conceptions that would exclude Religion, or even Art, from the Life of Reason, Santayana addresses all the important manifestations of human life. Moreover, the Life of Reason covers history and political science and cannot be restricted just to science and technology.

Hence a commitment to the Life of Reason benefits from the three rings of Saturn. Ultimately, Reason leads to Love. For Spinoza, the intellectual love of God is one's highest achievement. Santayana likewise points to the necessity of Love, although not in such lofty terms.

\paragraph{Animal Faith}


\begin{quotex}
That there is a world, that there is a future, that things sought can be found, and things seen can be eaten.
\flright{\textit{Skepticism and Animal Faith}}
\end{quotex}


We are born in the midst of a world; there is no objective and detached standpoint from which we can view the world. Animal Life precedes the Life of Reason. The above quotation applies both to animals and humans. Moreover, there are unexamined assumptions that Santayana makes explicit. For example, a belief in the existence of the material world and the reliability of the senses.

All that makes sense in itself. He is making explicit one of the inner wits of the Medievals, the one that C S Lewis translates as ``estimation''. Yet as much as Santayana insists that he is a materialist, he can be read, or perhaps creatively misread, as an idealist as I am wont to do. Why stop at our assumptions about the material world, since we can continue to play his game.

The philosopher has faith in the rationality and intelligibility of the world and the trustworthiness of human thinking, otherwise he would not bother with detailed intellectual studies.

Can we leave out the religious mind? If there is reason in religion, then religious faith is just as certain and ``a priori'', that is, prior to any possible proof. Do we doubt the perversity of the human condition, the sense of a fall from a higher state, the desire for salvation and ultimate liberation?

\paragraph{Realms of Being}


After elucidating the Life of Reason and our animal faith, Santayana then turns to existence, which comprises four realms of Being (the Realm of Truth adds nothing to this discussion). He believed he was creating a new system of philosophy, but all he did was rearrange traditional metaphysics.

\textbf{Essence} is the first realm. What we know are essences. This is what he owes to Plato; Schopenhauer likewise considers essences to be fundamental. Arguists will emphasizes the differences. Yet if we have a faith in the objective reality of the world, then such differences must needs be verbal, not real.

\textbf{Matter} is fundamental or primordial. It is Being itself. However, its various arrangements do not speak for themselves, but only through essences.

\textbf{Spirit} is the third realm. Santayana seems to mean consciousness. So we see the unavoidable issue. Matter and Spirit are variations on a theme, whether the attributes of Extensions and Thought in Spinoza, or the two Cartesian substances. Santayana presumes that matter creates spirit, or consciousness. Hence, a human person arises, lives a life for no particular reason, and then passes away into oblivion. Santayana denies that Spirit has agency; hence it can only be the passive observer of the material flux.

The esoterist can gain insights from the rings of Saturn. The Realm of Essence corresponds to the Divine Intellect (or Buddhi), the Realm of Matter to the Gross Body, and Spirit, or the real I, to the human state.


\flrightit{Posted on 2021-04-30 by Cologero}

